\section{Contenedorización con Docker Compose}
\label{sec:docker}

La principal ventaja de usar Docker~\cite{merkel2014docker,docker_compose_docs} es que
el sistema funciona igual en cualquier máquina sin necesidad de instalar nada: ni
GStreamer, ni Python, ni ninguna librería adicional. Con un único comando
(\texttt{make docker-up}) el sistema arranca completo en cuestión de segundos. Todo el
despliegue se define en dos ficheros: el \texttt{Dockerfile}, que describe cómo se
construye la imagen del sistema, y el \texttt{docker-compose.yml}, que declara los
servicios y cómo se relacionan entre sí. Ambos ficheros se reproducen íntegramente en
la sección~C6 del Anexo~C.

A partir del \texttt{Dockerfile} se genera una única imagen personalizada, basada en
Ubuntu~22.04~\ac{LTS}, que incluye todo lo necesario para ejecutar el sistema:
GStreamer~1.0 con sus \textit{plugins}, FFmpeg~\cite{ffmpeg_project}, las librerías
Python necesarias y \texttt{pika} para la integración con RabbitMQ. Esta misma imagen
sirve de base para todos los servicios: voctocore ejecuta el núcleo de mezcla, los
contenedores de cámara inyectan la señal de vídeo con FFmpeg, y el servicio de
telemetría registra los eventos. La imagen se construye automáticamente la primera vez
que se arranca el sistema. RabbitMQ es el único servicio que utiliza una imagen oficial
externa.

El principal problema técnico que hubo que resolver fue el aislamiento de red. En
Docker, cada contenedor está aislado y cuando se conecta a \texttt{localhost} solo se
comunica consigo mismo. Esto suponía un problema porque las cámaras necesitan enviar
la señal a voctocore a través de \texttt{localhost}. La solución fue hacer que todos
los contenedores compartan la red de voctocore mediante la opción
\texttt{network\_mode: "service:voctocore"}, de forma que todos los procesos se ven
entre sí como si estuvieran en la misma máquina, sin modificar ninguna línea del
código original.

Voctogui es el único componente que corre directamente en el equipo, sin contenedor,
porque necesita acceder al servidor gráfico del sistema para mostrar la interfaz. Se
conecta a voctocore a través del puerto~9999 y comparte los recursos visuales como
logos e imágenes de fondo mediante un volumen Docker, por lo que se pueden modificar
sin tener que reconstruir la imagen (figura~\ref{fig:escenario_docker}).

\begin{figure}[H]
  \centering
  \includegraphics[width=\textwidth]{figuras/escenario_docker.png}
  \caption{Escenario de despliegue con Docker Compose.}
  \label{fig:escenario_docker}
\end{figure}

Durante las pruebas se detectó que algunos contenedores intentaban conectarse a
voctocore antes de que estuviera listo, lo que causaba fallos en el arranque. Para
resolverlo se combinaron dos medidas: Docker reinicia automáticamente cualquier
contenedor que falle, y antes de lanzar cada servicio se verifica que el puerto~9999
está disponible. El arranque sigue así una secuencia ordenada: primero se levanta
voctocore, el resto de contenedores esperan a que el núcleo esté operativo, y
finalmente se lanza voctogui de forma nativa junto con un monitor de programa en el
puerto~15000. Al cerrar la sesión, todos los procesos se detienen de forma coordinada.
